\section{Projektudførelse}
\subsection{Oprettelse af projekt}
For at komme i gang med projektet oprettede gruppen først selve Unreal Engine-projektet. Projektet blev initialt oprettet som et C++-projekt baseret på en First Person-template, da gruppen ønskede, at spillet skulle opleves fra et førstepersonsperspektiv. Dette valg blev truffet tidligt i processen for at sikre, at gameplay-oplevelsen blev defineret allerede fra starten, hvilket senere har haft betydning for designvalg og udvikling af spilmekanikker.

Efter oprettelsen af projektet blev der opsat et fælles GitHub-repository for at understøtte samarbejdet i gruppen. Dette gjorde det muligt at arbejde parallelt på projektet, samtidig med at der blev skabt en struktureret versionstyring af filer og ændringer. Brugen af pull requests og code reviews blev implementeret for at sikre kvalitet i koden og for at give mulighed for faglig sparring mellem gruppemedlemmerne. Denne arbejdsform viste sig at være vigtig for at undgå konflikter i koden og for at sikre, at ændringer blev gennemgået og forstået af flere i gruppen, før de blev integreret i hovedprojektet.

Tidligt i udviklingsfasen blev der eksperimenteret med Unreal Engine, særligt i forhold til samspillet mellem Blueprints og C++. Gruppen undersøgte, hvordan Blueprints kunne anvendes til hurtig prototyping af gameplay-elementer, mens C++ blev brugt til mere komplekse og performancekritiske systemer. Denne kombination viste sig at være central for projektets udvikling, da den gav en god balance mellem fleksibilitet og kontrol.

I begyndelsen af projektet bar arbejdet præg af eksperimentering og læring, hvor gruppen brugte tid på at opbygge forståelse for engine’ens struktur og arbejdsmetoder. Dette betød, at fremdriften i starten var relativt langsom, men til gengæld lagde det et solidt fundament for den videre udvikling. Efterhånden som gruppen fik større teknisk forståelse, blev arbejdsprocessen mere effektiv, og udviklingen af spillets funktioner kunne ske mere struktureret og målrettet.

\newpage
\subsection{Game Sessions}
den første del af systemet som blev udviklet var Faktisk Online Session systmet, dette blev gjordt, for at sikre at alle efterfølgende systemer kunne testes i en online kontekst, og for at sikre at vi havde en stabil base at bygge videre på. Dette system blev udviklet ved hjælp af Unreal Engine's Online Subsystem, som giver en række værktøjer og funktioner til at håndtere online multiplayer-funktionalitet. Implementeringen omfattede oprettelse og håndtering af sessions, matchmaking, og netværkskommunikation mellem klienter og servere. Dette var en vigtig milepæl i projektet, da det lagde grundlaget for alle de multiplayer-aspekter, der senere skulle implementeres i spillet.

i praksis blev implementationen af online sessions først oprettet som et system hvor spilleren kunne oprette en session, og spiller der joinede ville altid senere 

\newpage
\subsection{TeamSelection}


\newpage
\subsection{VoiceChat}
//tænker Victor dit kan være her


\newpage
\subsection{Timer}
//tænker Jonathan dit kan være her


\newpage
\subsection{Class Selection}


\newpage
\subsection{Speceille Evner}
//tænker Jonathan dit kan være her


\newpage
\subsection{nøgler og kisten}


\newpage
\subsection{Fange Mechanic}


\newpage
\subsection{befri fængslet}


\newpage
\subsection{Win Conditions}


\newpage
\subsection{map opsætning}


\newpage
\subsection{AI NPC'er}
